\documentclass[11pt,a4paper]{article}

\usepackage[utf8]{inputenc}
\usepackage[T1]{fontenc}
\usepackage{lmodern}
\usepackage[margin=1in]{geometry}
\usepackage{amsmath,amssymb,amsthm}
\usepackage{hyperref}
\usepackage{xcolor}
\usepackage{listings}
\usepackage{booktabs}
\usepackage{enumitem}
\usepackage{float}
\usepackage{tikz}
\usetikzlibrary{arrows.meta,positioning,shapes.geometric,calc}
\usepackage{fancyhdr}
\usepackage{titlesec}
\usepackage{tabularx}

\definecolor{modblue}{HTML}{1E3A5F}
\definecolor{modcyan}{HTML}{0EA5E9}
\definecolor{codebg}{HTML}{F8F9FA}
\definecolor{codeframe}{HTML}{DEE2E6}

\hypersetup{
    colorlinks=true,
    linkcolor=modblue,
    citecolor=modblue,
    urlcolor=modcyan,
    pdftitle={Scaling Off-Chain Open-Source Management with Merkle Trees},
    pdfauthor={MOD Protocol},
}

\lstdefinestyle{solidity}{
    backgroundcolor=\color{codebg},
    frame=single,
    rulecolor=\color{codeframe},
    basicstyle=\ttfamily\small,
    keywordstyle=\color{modblue}\bfseries,
    commentstyle=\color{gray},
    stringstyle=\color{teal},
    breaklines=true,
    numbers=left,
    numberstyle=\tiny\color{gray},
    numbersep=8pt,
    tabsize=4,
    showstringspaces=false,
    morekeywords={function,external,internal,public,private,view,pure,returns,
                  require,mapping,address,uint256,bytes32,bool,memory,calldata,
                  struct,event,emit,modifier,contract,interface,pragma,solidity},
}

\lstdefinestyle{python}{
    backgroundcolor=\color{codebg},
    frame=single,
    rulecolor=\color{codeframe},
    basicstyle=\ttfamily\small,
    keywordstyle=\color{modblue}\bfseries,
    commentstyle=\color{gray},
    stringstyle=\color{teal},
    breaklines=true,
    numbers=left,
    numberstyle=\tiny\color{gray},
    numbersep=8pt,
    tabsize=4,
    showstringspaces=false,
    language=Python,
}

\pagestyle{fancy}
\fancyhf{}
\fancyhead[L]{\small\textcolor{gray}{MOD Protocol -- Merkle-Tree Module Registry}}
\fancyhead[R]{\small\textcolor{gray}{v0.1 -- May 2026}}
\fancyfoot[C]{\thepage}
\renewcommand{\headrulewidth}{0.4pt}

\titleformat{\section}{\Large\bfseries\color{modblue}}{\thesection}{1em}{}
\titleformat{\subsection}{\large\bfseries\color{modblue!80}}{\thesubsection}{1em}{}

\newtheorem{definition}{Definition}[section]
\newtheorem{property}{Property}[section]
\newtheorem{claim}{Claim}[section]

\begin{document}

\begin{titlepage}
\centering
\vspace*{3cm}

{\Huge\bfseries\textcolor{modblue}{Scaling Off-Chain\\Open-Source Management}\par}
\vspace{0.6cm}
{\Large\textcolor{gray}{Storing the tree, not the leaves:\\a Merkle-root registry for the MOD orbit ecosystem}\par}

\vspace{2cm}
{\large Version 0.1\par}
{\large May 2026\par}

\vspace{3cm}

\begin{abstract}
\noindent
The MOD protocol today registers each orbit module as an independent on-chain row keyed by an IPFS CID. With 200+ modules and projected growth into the tens of thousands as autonomous agents publish their own, the linear per-module storage and gas cost becomes the binding constraint on ecosystem scale. We propose replacing the per-module hash table with a single 32-byte Merkle root anchored on-chain. The full tree---containing module metadata, code CIDs, owner addresses, prices, and version history---lives entirely off-chain in IPFS-backed mirrors. Authenticity is preserved by Merkle inclusion proofs; updates are batched into a single root-rotation transaction; reads are free. We show that this design reduces per-module on-chain storage from $\approx 5$ slots ($\approx 100{,}000$ gas amortised) to $O(1/N)$, reduces a publish operation from $\approx 80{,}000$ gas to a marginal cost approaching the cost of a single SSTORE shared across $N$ updates, and makes the registry compatible with permissionless off-chain indexers without trust assumptions. The trade-off is a one-block latency between local commit and global root rotation, which is acceptable for a developer-tools registry.
\end{abstract}

\vfill
\end{titlepage}

\tableofcontents
\newpage

\section{Motivation}

\subsection{Current design}

The current \texttt{Registry.sol} stores one struct per module:

\begin{lstlisting}[style=solidity]
struct Module {
    address owner;
    string  name;
    string  data;   // IPFS CID
    uint256 modId;
}
mapping(uint256 => Module) public modules;
mapping(address => uint256[]) public userMods;
\end{lstlisting}

Each \texttt{registerMod} call writes:
\begin{itemize}
    \item the struct (typically 4--5 SSTOREs after string encoding overhead),
    \item an append to \texttt{userMods},
    \item a monotonic counter bump.
\end{itemize}

At 80{,}000--110{,}000 gas per registration on Base, and re-registration on every CID update, the cost is dominated by storage growth in the registry contract itself.

\subsection{The scaling wall}

Three trends make the per-row design untenable:

\begin{enumerate}
    \item \textbf{Agentic publishing.} The \texttt{agent} module already composes new modules autonomously; a 24-hour run can produce tens of new mods.
    \item \textbf{Versioning.} Most production modules update weekly. The registry currently has no compression: each version is a fresh write.
    \item \textbf{Permissionless mirrors.} Indexers that want to expose the registry over HTTP must replay every \texttt{ModRegistered} event since genesis, which is a quadratic cost in time and state.
\end{enumerate}

\subsection{Goal}

Compress the on-chain footprint to a single root, while preserving:

\begin{itemize}
    \item authenticity (no off-chain actor can lie about what's registered),
    \item revocability (owners can delete or transfer modules),
    \item ownership (only the rightful owner can update an entry),
    \item discoverability (clients can enumerate without trust).
\end{itemize}

\section{Construction}

\subsection{The off-chain tree}

\begin{definition}[Module record]
A module record $r$ is a tuple
\[
r = (\text{name},\ \text{owner},\ \text{cid},\ \text{price},\ \text{version},\ \text{updatedAt})
\]
serialised with canonical JSON (RFC 8785) and hashed with \texttt{keccak256}:
\[
\ell(r) = \mathrm{keccak256}(\mathrm{cjson}(r)).
\]
\end{definition}

\begin{definition}[Module tree]
Given $N$ module records $\{r_1, \ldots, r_N\}$, the module tree $\mathcal{T}$ is a binary Merkle tree over the multiset $\{\ell(r_1), \ldots, \ell(r_N)\}$ sorted by $\ell$ (sorted-pair construction; identical to OpenZeppelin's \texttt{MerkleProof} library). The root is
\[
R = \mathrm{root}(\mathcal{T}).
\]
\end{definition}

\subsection{On-chain anchor}

\begin{lstlisting}[style=solidity]
contract TreeRegistry {
    bytes32 public root;          // current Merkle root
    uint256 public epoch;         // rotation counter
    address public publisher;     // authorised root publisher (later DAO)
    mapping(bytes32 => bool) public revoked;  // tombstones

    event Rotated(uint256 indexed epoch, bytes32 indexed root, bytes32 cid);
    event Revoked(bytes32 indexed leaf);

    function rotate(bytes32 newRoot, bytes32 manifestCid) external {
        require(msg.sender == publisher);
        root = newRoot;
        epoch++;
        emit Rotated(epoch, newRoot, manifestCid);
    }

    function verify(bytes32 leaf, bytes32[] calldata proof)
        external view returns (bool)
    {
        if (revoked[leaf]) return false;
        return MerkleProof.verifyCalldata(proof, root, leaf);
    }

    function revoke(bytes32 leaf, bytes32[] calldata proof) external {
        // Only the owner encoded in the leaf can revoke.
        // Owner check performed off-chain by reconstructing the record
        // and verifying that msg.sender matches r.owner.
        // ...
        revoked[leaf] = true;
        emit Revoked(leaf);
    }
}
\end{lstlisting}

The on-chain state is exactly:
\begin{itemize}
    \item one \texttt{bytes32} root,
    \item one \texttt{uint256} epoch,
    \item one \texttt{address} publisher,
    \item a sparse \texttt{revoked} map (only paid for by the revoker, only for actively revoked entries).
\end{itemize}

\subsection{The publisher and the manifest}

The publisher is the actor authorised to call \texttt{rotate}. In Phase 1 it is the protocol multisig; in Phase 2 it is a threshold of validators sourced from \texttt{StakeTime}.

On every rotation the publisher uploads the full tree to IPFS as a \emph{manifest}:

\begin{lstlisting}[basicstyle=\ttfamily\small,frame=single,backgroundcolor=\color{codebg}]
{
  "epoch": 412,
  "root": "0x9f4c...e2",
  "prev": "0x2cb1...0a",
  "records": [
    { "name": "agent",     "owner": "0x...", "cid": "Qm...", ... },
    { "name": "bloctime",  "owner": "0x...", "cid": "Qm...", ... },
    ...
  ],
  "tree": { "depth": 18, "nodes": "..." }
}
\end{lstlisting}

The manifest CID is emitted in the \texttt{Rotated} event so that any indexer with an IPFS gateway can pull the new tree and verify its root against the on-chain value.

\subsection{Update flow}

\begin{enumerate}
    \item Developer signs a record update with their wallet.
    \item Signed update is broadcast to a public mempool (a Caddy-fronted FastAPI endpoint at \texttt{/api/whitepaper/tree/submit}).
    \item Publisher (or any node) aggregates updates between epochs.
    \item Publisher recomputes the tree, uploads the manifest, calls \texttt{rotate(newRoot, manifestCid)}.
    \item Clients verify their relevant leaves against the new root via Merkle proofs.
\end{enumerate}

A single \texttt{rotate} transaction subsumes an arbitrary number of registrations and updates. The cost per individual update is $O(1/N)$ where $N$ is the rotation batch size.

\section{Gas and Storage Analysis}

\subsection{Per-operation gas}

\begin{table}[H]
\centering
\begin{tabularx}{\textwidth}{lXX}
\toprule
\textbf{Operation} & \textbf{Current Registry} & \textbf{Tree Registry} \\
\midrule
Register module          & 80{,}000--110{,}000 & $\approx 30{,}000 / N$ (amortised) \\
Update CID               & 30{,}000--50{,}000 & $\approx 30{,}000 / N$ (amortised) \\
Transfer ownership       & 25{,}000 & $\approx 30{,}000 / N$ \\
Revoke / delete          & 15{,}000 (SSTORE clear) & 25{,}000 (single SSTORE in \texttt{revoked} map) \\
Read by name             & free (view) & free (Merkle verify off-chain) \\
Enumerate all modules    & $O(N)$ event replay & 1 IPFS fetch of manifest \\
\bottomrule
\end{tabularx}
\caption{Gas comparison. $N$ is the rotation batch size (e.g.\ 100 updates per rotation).}
\end{table}

For $N = 100$, a rotation cycle costs $\approx 30{,}000 / 100 = 300$ gas per update---two orders of magnitude below the current floor.

\subsection{On-chain storage growth}

\begin{property}[Constant footprint]
The contract storage of \texttt{TreeRegistry} grows only with the number of actively revoked modules (a tiny minority). Compared to the linear $O(N)$ growth of the per-row registry, on-chain state is effectively constant.
\end{property}

\begin{property}[Sublinear emission cost]
A \texttt{Rotated} event is fixed-size regardless of batch contents. The \texttt{ModRegistered} event stream of the current design is $O(N)$ and quadratic to fully replay.
\end{property}

\section{Verifiability and Trust}

\subsection{What a client needs}

To verify that module \texttt{X} is registered to owner \texttt{O} with CID \texttt{C}:

\begin{enumerate}
    \item Fetch the current \texttt{root} from \texttt{TreeRegistry} (1 \texttt{eth\_call}).
    \item Fetch the manifest CID from the latest \texttt{Rotated} event (1 log filter).
    \item Pull the manifest from any IPFS gateway.
    \item Reconstruct the record locally and compute $\ell(r)$.
    \item Compute the Merkle proof against the on-chain root.
    \item (Optional) check that $\ell(r)$ is not in \texttt{revoked}.
\end{enumerate}

The client does \emph{not} need to trust the publisher: if the publisher signs a root that does not match the manifest, mismatched proofs reveal it instantly.

\begin{claim}[No fraudulent record can be proven]
\label{claim:no-fraud}
For any record $r' \notin \mathcal{T}$, no Merkle proof against $\mathrm{root}(\mathcal{T})$ verifies, under the standard collision-resistance assumption on \texttt{keccak256}.
\end{claim}

\subsection{Censorship}

The publisher can refuse to include a submitted update. We mitigate this two ways:

\begin{enumerate}
    \item \textbf{Forced inclusion.} Any user may pay gas to call \texttt{forceInclude(record, signature)} on a fallback contract that bypasses the publisher and triggers a special rotation. This restores the worst-case behaviour to today's per-row cost---but only when censored, not by default.
    \item \textbf{Publisher rotation.} If the publisher is a validator threshold (Phase 2), bribing $k$-of-$n$ validators is required to censor. Slashing is on the line.
\end{enumerate}

\section{Off-Chain Storage Architecture}

\subsection{Storage tiers}

\begin{table}[H]
\centering
\begin{tabularx}{\textwidth}{lX}
\toprule
\textbf{Tier} & \textbf{Contents} \\
\midrule
On-chain (1 slot) & Current root, epoch, publisher address \\
IPFS (manifests) & Full module records + tree structure, addressed by CID \\
Filecoin (long-term) & Manifest pinning with economic guarantees \\
Local cache & Compact LMDB of (leaf $\to$ proof) for each subscribed indexer \\
\bottomrule
\end{tabularx}
\caption{Storage tiering. The smaller the tier, the higher the trust.}
\end{table}

\subsection{Replication and availability}

A manifest is small (estimated $\approx 50$ bytes/record + tree overhead, so a 100k-module ecosystem is $\approx 5$ MB). Every running orbit node can pin every historical manifest at negligible cost. The protocol does not require a single trusted gateway; any IPFS or HTTP mirror suffices because authenticity is verified against the on-chain root.

\subsection{Module API surface}

This whitepaper ships as a runnable orbit module (\texttt{mod/orbit/whitepaper/}) with both a static document app and a working reference implementation of the tree:

\begin{lstlisting}[style=python]
m.fn('whitepaper/tree_build')(records)          # build tree + manifest
m.fn('whitepaper/tree_root')()                  # current root
m.fn('whitepaper/tree_proof')(name='agent')     # Merkle proof for a leaf
m.fn('whitepaper/tree_verify')(leaf, proof)     # client-side verify
\end{lstlisting}

\section{Migration Path}

\begin{enumerate}
    \item \textbf{Shadow mode.} \texttt{TreeRegistry} deploys alongside the existing per-row \texttt{Registry}. The publisher mirrors all \texttt{Registry} events into the tree. Both registries serve simultaneously.
    \item \textbf{Indexer flip.} Off-chain indexers (the \texttt{api} module, the polymarket frontend, ecosystem block explorers) cut their reads over to the tree-based view. Writes still go to the legacy registry.
    \item \textbf{Write deprecation.} New \texttt{registerMod} calls are routed through a meta-aggregator that submits to the tree pool. The legacy contract is frozen but its state is mirrored permanently into the tree's genesis manifest.
    \item \textbf{Sunset.} The legacy registry stays readable; no new writes are accepted. Eventually it is marked \texttt{setOwnerless} and forgotten.
\end{enumerate}

\section{Comparison to Related Work}

\begin{itemize}
    \item \textbf{Ethereum Name Service (Namewrapper).} ENS stores per-name records on-chain. Our design is closer to \emph{rollup} commit roots applied to a developer registry.
    \item \textbf{Optimistic rollups.} We borrow the publish-root-then-prove-fraud pattern, but without fraud proofs: there is no state transition to dispute, only a static set membership. This makes the model significantly simpler.
    \item \textbf{IPNS / IPFS DNSLink.} Both lack an on-chain anchor, so they cannot resolve disputes about which manifest is canonical. Our root + epoch counter solves that.
    \item \textbf{Sparse Merkle Trees (SMT).} An SMT would give us inclusion + non-inclusion proofs for free at the cost of larger proofs. For an opt-in publisher model where revocation is a tombstone bit, we do not need non-inclusion proofs; the sorted-leaf binary tree is sufficient and cheaper.
\end{itemize}

\section{Conclusion}

Storing the tree instead of the leaves shifts the cost curve from $O(N)$ to $O(1)$ on chain, eliminates per-module gas as a publishing tax, and makes the registry trivially mirrorable. It preserves authenticity through Merkle proofs and ownership through signed updates, and degrades gracefully under censorship via forced inclusion. The resulting system is the registry layer the MOD orbit ecosystem needs to grow from hundreds of modules to millions, including those published autonomously by agents.

\appendix

\section{Reference Implementation}

The module ships a Python reference implementation at \texttt{mod/orbit/whitepaper/mod.py} exposing:

\begin{lstlisting}[style=python]
class Mod:
    def tree_build(self, records: list[dict]) -> dict: ...
    def tree_root(self) -> str: ...
    def tree_proof(self, name: str) -> list[str]: ...
    def tree_verify(self, leaf: str, proof: list[str]) -> bool: ...
\end{lstlisting}

\section{Notation}

\begin{tabularx}{\textwidth}{lX}
\toprule
$N$ & Total number of modules in the registry. \\
$r$ & A module record (canonical JSON). \\
$\ell(r)$ & The leaf hash $\mathrm{keccak256}(\mathrm{cjson}(r))$. \\
$\mathcal{T}$ & The Merkle tree over all leaves. \\
$R$ & The root of $\mathcal{T}$, anchored on-chain. \\
$E$ & The epoch counter. \\
\bottomrule
\end{tabularx}

\end{document}
